\documentclass{tutodoc}

\usepackage{../preamble.cfg}


\begin{document}

\section{Where do the color palettes come from?}
\label{projects-used}

\subsection{Resources used}

\thisproj\ includes some original creations, but most color palettes are derived from the projects listed below.
%
% ---------------------------------------- %
% -- JUST TRANSLATE THE ''XTRA'' TEXTS. -- %
% ---------------------------------------- %
%
% \def\fromcolormaps{palettes are extracted from \colormaps\ project}
\def\frompalettable%
       {palettes are extracted from \palettable\ project}
%
\begin{multicols}{2}
\begin{itemize}
  \foreach \techno/\xtra in {%
    % -- A -- %
    \asymptote / %
        {is used, but currently offers nothing beyond \matplotlib\ (despite different implementations)},%
    % -- C -- %
    \carbonplan  / %
        {provides palettes designed for data visualization.},
    \cartocolors / \frompalettable,%
    \cmasher     / %
        {provides palettes designed for data visualization.},
    \cmocean     / \frompalettable,%
    \colorbrewer / %
        {provides professional color palettes for mapping and data visualization},%
    % \colorcet    / \fromcolormaps, %
    % -- L -- %
    \lightbartlein / \frompalettable,%
    % -- M -- %
    \matplotlib / %
        {compiles color maps from diverse projects}, %
    \mycarta    / \frompalettable,%
    % -- P -- %
    \plotly / \frompalettable,%
    % -- S -- %
    \scicolmaps / %
        {provides palettes designed for colorblind accessibility}, %
    % \sciviz    / \fromcolormaps, %
    % -- T -- %
    \tableau / \frompalettable, %
    % -- W -- %
    \wesanderson / \frompalettable %
  }{%
    \item \techno\xtra.
  }
\end{itemize}
\end{multicols}


% -------------------- %


\begin{tdocnote}
    The following versions were used to build the current \thisproj\ release.
    
    \smallskip
    
    \input{../../../common/rescr/last-versions.latex}
\end{tdocnote}

\end{document}
